\section{Encoded frames of guides}
\label{sec:peps_guide_frames}
The statements below set up the encoded frames used in the proof of
Proposition~7.1 of \cite{OpenAI2026PolynomialPEPS}, September 24, 2026,
\texttt{06-geometry.tex}, lines 62--125 and 669--672: the raw owners are read
from a guide at the sampled sites, the holes sit at given centers with the
patches supplied by Proposition~4.1 of the same source, and the tags are held by
a canonical incident label. Proposition~4.1 itself is stated here and enters as
a hypothesis; its proof is not part of this section.

We keep the setting of Section~\ref{sec:peps_encoded_frames}. For square cylinder
data with radii $r_1,\ldots,r_n$ and vectors $v_{j\ell}$, write
\[
    \Pi=\sum_{j=1}^{n}\sum_{\ell=1}^{d_j}
    \bigl(\ket{v_{j\ell}}\bra{v_{j\ell}}\bigr)_{D_j}\otimes\id
\]
for the sum of the cylinder projectors, without the convention for an empty
outer sample.

\begin{definition}[Statement of the adaptive patch constraint]
    \label{def:peps_adaptive_patch_constraint}
    \lean{TNLean.PEPS.EncodedFrame.SquarePatch.cylSum,
        TNLean.PEPS.EncodedFrame.HasAdaptivePatch,
        TNLean.PEPS.EncodedFrame.AdaptivePatchSupply,
        TNLean.PEPS.Approximation.latticePos,
        TNLean.PEPS.Approximation.AdaptivePatchConstraint}
    \leanok
    \uses{def:peps_frame_square_patch}
    Let $\Lambda_L=\{1,\ldots,L\}^2$ with its sites at their integer positions.
    The adaptive patch constraint is the following assertion. Fix $q\ge2$,
    $J,\Delta>0$, $B_0>0$ and $p>0$. There are constants $C,c>0$ such that for
    every $L\ge2$, every square Hamiltonian with local bounds $J$ and every
    unit ground vector $\Omega$ with gap $\Delta$, every center $z\in\R^2$ and
    every $0<u\le B_0\log L$ admit square cylinder data at $(z,u)$ with
    \[
        \sum_{j}d_j\le CL^c,
        \qquad
        \lVert(\id-\Pi)\Omega\rVert\le L^{-p}.
    \]
    This is \cite[Proposition~4.1]{OpenAI2026PolynomialPEPS},
    \texttt{03-patches.tex}, lines 26--49, with the parameters of the main
    theorem fixed as in \texttt{01-preliminaries.tex}, lines 3--10. The sum
    $\Pi$ is an orthogonal projection supported on the sample of the square of
    radius $2u$ because the cylinder projectors have orthogonal ranges.
\end{definition}

\begin{lemma}[Supplied patches]
    \label{lem:peps_adaptive_patch_supply}
    \lean{TNLean.PEPS.EncodedFrame.SquarePatch.proj_eq_cylSum,
        TNLean.PEPS.EncodedFrame.SquarePatch.norm_act_proj_sub_le,
        TNLean.PEPS.EncodedFrame.SquarePatch.card_tag_le_max,
        TNLean.PEPS.EncodedFrame.SquarePatch.translate,
        TNLean.PEPS.EncodedFrame.HasAdaptivePatch.translate,
        TNLean.PEPS.EncodedFrame.patchAt,
        TNLean.PEPS.EncodedFrame.norm_act_proj_patchAt_sub_le,
        TNLean.PEPS.EncodedFrame.card_tag_patchAt_le,
        TNLean.PEPS.Approximation.cellCenterPos,
        TNLean.PEPS.Approximation.AdaptivePatchSupply.cellCenterPos,
        TNLean.PEPS.Approximation.AdaptivePatchConstraint.supply_cellCenterPos}
    \leanok
    \uses{def:peps_adaptive_patch_constraint, def:peps_frame_hole_encoder}
    If the outer sample of square cylinder data is nonempty, its projector
    $P_{c,h}$ is $\Pi$; in all cases
    $\lVert(\id-P_{c,h})\Omega\rVert\le\lVert(\id-\Pi)\Omega\rVert$ and the
    number of tags is at most $\max\{1,\sum_jd_j\}$. Translating the site
    positions and the center by one vector does not change the samples, so the
    conclusion of the adaptive patch constraint is translation invariant. Hence,
    if the adaptive patch constraint holds, then with the same constants it holds
    on the sites placed at the cell centers $(i+\tfrac12,j+\tfrac12)$,
    $0\le i,j<L$, as in \texttt{06-geometry.tex}, lines 62--63. Fixing one
    supplied patch for each center and scale, the projector of the patch at
    $(z,u)$ has error at most $L^{-p}$ on $\Omega$ and at most $\max\{1,CL^c\}$
    tags. The patch depends only on the center and the scale, so a given center
    and radius receive the same patch on every sheet, as required at
    \texttt{06-geometry.tex}, lines 114--118.
\end{lemma}

\begin{proof}\leanok
    With a nonempty outer sample the tags are the pairs $(j,\ell)$; with an
    empty one the projector is the identity. A translation preserves sup-norm
    distances.
\end{proof}

\begin{definition}[Canonical tag holder and the frame of a guide]
    \label{def:peps_guide_frame}
    \lean{TNLean.PEPS.Approximation.tagHolder,
        TNLean.PEPS.Approximation.guideFrame}
    \leanok
    \uses{def:peps_dyadic_incident_labels, def:peps_frame}
    Fix a ranking of the labels by natural numbers and a placeholder label. The
    canonical tag holder of a point $c$ of a guide is the incident label at $c$
    other than the placeholder of least rank. Given a guide read at the sites, an
    assignment of parties to labels, a finite set $V$ of centers and a square
    patch at each center, with pairwise disjoint outer samples, the frame of the
    guide has the raw owner of a site equal to the party of the label read at its
    position, and one hole at each center of $V$, carrying its patch, with tag
    held by the party of the canonical tag holder. This is the frame described
    in \cite[Section~7]{OpenAI2026PolynomialPEPS}, \texttt{06-geometry.tex},
    lines 69--80 and 109--125.
\end{definition}

\begin{lemma}[Canonical tag holders]
    \label{lem:peps_guide_frame_tags}
    \lean{TNLean.PEPS.Approximation.isTagHolder_tagHolder,
        TNLean.PEPS.Approximation.tagHolder_ne,
        TNLean.PEPS.Approximation.rank_tagHolder_le,
        TNLean.PEPS.Approximation.tagHolder_congr}
    \leanok
    \uses{def:peps_guide_frame, lem:peps_dyadic_tags}
    At a true vertex the canonical tag holder is an incident label other than the
    placeholder, and it has least rank among such labels. Guides with the same
    incident labels at a point have the same canonical tag holder there, so the
    tag of a hole accompanies it through every change that keeps the incident
    labels, as at \texttt{06-geometry.tex}, lines 123--125.
\end{lemma}

\begin{proof}\leanok
    A true vertex has an incident label other than the placeholder, so the set
    of candidates is nonempty and has an element of least rank. The set of
    candidates depends only on the incident labels.
\end{proof}

\begin{theorem}[Reference vectors of guide frames]
    \label{thm:peps_guide_frame_norm}
    \lean{TNLean.PEPS.Approximation.disjoint_outer_of_disjoint_closedBall,
        TNLean.PEPS.Approximation.pairwise_disjoint_outer_patchAt,
        TNLean.PEPS.Approximation.mem_holes_guideFrame,
        TNLean.PEPS.Approximation.length_holes_guideFrame,
        TNLean.PEPS.Approximation.one_sub_mul_le_norm_refVec_guideFrame,
        TNLean.PEPS.Approximation.one_sub_mul_le_norm_refVec_guideFrame_patchAt,
        TNLean.PEPS.Approximation.outer_eq_empty_of_isCellCenter,
        TNLean.PEPS.Approximation.guideFrame_unitLevel}
    \leanok
    \uses{def:peps_guide_frame, thm:peps_frame_norm, lem:peps_adaptive_patch_supply}
    Patches whose closed outer squares are disjoint in the plane have disjoint
    outer samples. Let $\lVert\Omega\rVert=1$. If every patch at a center of $V$
    has projector error at most $\varepsilon$ on $\Omega$, the reference vector of
    the frame of a guide satisfies
    $1-\lvert V\rvert\varepsilon\le\lVert\Omega_{\mathcal F}\rVert\le1$; with
    supplied patches one may take $\varepsilon=L^{-p}$. If the sites lie at cell
    centers, every center of $V$ is a corner of a unit block, and every outer
    radius is less than $\tfrac12$, then every hole has empty physical outer
    sample, the encoding is an isometry and
    $\lVert\Omega_{\mathcal F}\rVert=\lVert\Omega\rVert$. This is the last level
    $n=1$ of \texttt{06-geometry.tex}, lines 669--672, where
    $2h_1\le2\epsilon_0<\tfrac12$.
\end{theorem}

\begin{proof}\leanok
    The first claim is the definition of the outer samples. The bound is
    Theorem~\ref{thm:peps_frame_norm}, since the frame has $\lvert V\rvert$ holes.
    A cell center is at sup distance at least $\tfrac12$ from every integer
    point, so the outer samples at $n=1$ are empty and every hole projector is
    the identity.
\end{proof}
